% ==============================================================================
% styles/tables.tex
% ==============================================================================
% این فایل شامل ۲ استایل زیبا و متفاوت برای جداول با پشتیبانی کامل از راست‌چین (RTL) است.

% تنظیمات حرفه‌ای کپشن جداول
\captionsetup[table]{
    labelfont={bf,color=ThemeSecondary},
    textfont=small,
    skip=8pt,
    justification=centering
}

% ---------------------------------------------------------
% راهنمای استفاده از ستون‌ها و عرض جداول:
% 
% در این قالب از پکیج tabularray برای ساخت جداول بسیار حرفه‌ای استفاده شده است.
% 
% 1. ستون‌های معمولی: 
%    c (وسط‌چین)، r (راست‌چین)، l (چپ‌چین)
%
% 2. ستون با عرض ثابت (محتوای طولانی می‌شکند):
%    p{3cm}
%
% 3. جدول هم‌عرض صفحه (Page Width):
%    با قرار دادن ستون از نوع X، جدول کشیده شده و عرض صفحه را پر می‌کند.
% ---------------------------------------------------------

% دستور کمکی برای هدر جداول رنگی (برای سازگاری با نسخه‌های قبلی، در tabularray رنگ هدر خودکار است)
\newcommand{\THead}[1]{\textcolor{white}{\textbf{#1}}}

% =========================================================
% 1. جدول مینیمال (Minimal Table)
% =========================================================
\NewDocumentEnvironment{minimaltable}{ O{H} m m +b }{%
    \begin{table}[#1]
        \centering
        \caption{#2}
        \vspace{0.3em}
        
        \def\mybody{#4}%
        \begin{tblr}[expand=\mybody]{
            colspec = {#3},
            row{1} = {font=\bfseries},
            hline{1,Z} = {1pt, solid},
            hline{2} = {0.5pt, solid},
            stretch = 1.35
        }%
        \mybody
        \end{tblr}
    \end{table}
}{}

% =========================================================
% 2. جدول مدرن (Modern Table)
% =========================================================
\NewDocumentEnvironment{moderntable}{ O{H} m m +b }{%
    \begin{table}[#1]
        \centering
        \caption{#2}
        \vspace{0.4em}
        
        \begin{tcolorbox}[
            enhanced,
            arc=1.5mm,          % گوشه‌های گرد و حرفه‌ای
            boxrule=1.5pt,      % ضخامت کادر جدول
            colframe=ThemeSecondary,
            colback=white,
            boxsep=0pt,
            left=0pt, right=0pt, top=0pt, bottom=0pt,
            clip upper,
            hbox
        ]%
        \def\mybody{#4}%
        \begin{tblr}[expand=\mybody]{
            colspec = {#3},
            row{1} = {bg=ThemeSecondary, fg=white, font=\bfseries},
            row{even} = {bg=ThemeSecondary!8},
            stretch = 1.35
        }%
        \mybody
        \end{tblr}%
        \end{tcolorbox}
    \end{table}
}{}
